\subsection{Harris片的外区解：亲自验出$\Delta'$}
对于$B_x=B_0\tanh(y/a)$，设$s=y/a,K=ka$。有
\begin{equation}
 {B_x''\over B_x}=-{2\over a^2}\operatorname{sech}^2s,\qquad
 {\dd^2\tilde\psi\over\dd s^2}
 -[K^2-2\operatorname{sech}^2s]\tilde\psi=0.
\end{equation}
对$s>0$取$f=e^{-Ks}(K+\tanh s)$。令$T=\tanh s$、$Q=1-T^2$，
\begin{align}
 f'&=e^{-Ks}[Q-K(K+T)],\\
 f''&=e^{-Ks}[-2TQ-2KQ+K^2(K+T)]\\
 &=(K^2-2Q)f.
\end{align}
这就直接验证了解满足外区方程。又有$f(0)=K$、$f'(0)=1-K^2$，所以
$f'(0)/(af(0))=(1/K-K)/a$。取关于$y=0$偶对称的磁通扰动，负侧导数符号相反，故跳跃为两倍：
\begin{equation}
 \Delta'={2\over a}(K^{-1}-K).
\end{equation}
例如$ka=0.5$给$\Delta'a=3$；$ka=1$给零。$\Delta'>0$表示该经典平衡存在撕裂驱动，但正$\Delta'$本身不规定增长率、内层宽度或非线性岛数。

\subsection{学习内外匹配：先把电阻内层标度完整消元}
下面只为学习匹配方法，取不可压、无Hall、无电子惯性、低黏性的经典电阻内层。以片半厚$a$和相应$\VA$归一，定义
$\epsilon=\delta/a$、$K=ka$、$\hat\gamma=\gamma a/\VA$、$S_a=a\VA/D_\eta$、$\Delta=\Delta'a$。近共振面$B_x/B_0\simeq y/a$。线性涡量与磁通方程的量级为
\begin{equation}
 \hat\gamma{\tilde\phi\over\epsilon^2}
 \sim K\epsilon\,\tilde\psi'',\qquad
 \hat\gamma\tilde\psi\sim K\epsilon\tilde\phi+S_a^{-1}\tilde\psi''.
\end{equation}
双撇此处是对无量纲横坐标的导数。相位因子$i$不影响增长率量级，真正求本征函数时必须保留。

在常$\psi$内层中，磁通变化很小，但导数跳跃有限。跨越宽度$\epsilon$后有
$\tilde\psi''\epsilon\sim\Delta\tilde\psi$，因此
$\tilde\psi''\sim\Delta\tilde\psi/\epsilon$，不是$\tilde\psi/\epsilon^2$。
由三项平衡得
\begin{equation}
 \tilde\phi\sim{K\Delta\epsilon^2\over\hat\gamma}\tilde\psi,\qquad
 \hat\gamma\sim{S_a^{-1}\Delta\over\epsilon},\qquad
 \hat\gamma^2\sim K^2\Delta\epsilon^3.
\end{equation}
把第二式平方并代入第三式：
\begin{equation}
 {S_a^{-2}\Delta^2\over\epsilon^2}
 \sim K^2\Delta\epsilon^3
 \ \Rightarrow\ 
 \epsilon\sim S_a^{-2/5}\Delta^{1/5}K^{-2/5},\quad
 \hat\gamma\sim S_a^{-3/5}\Delta^{4/5}K^{2/5}.
\end{equation}
常$\psi$的自洽条件是内层磁通变化占比
$\Delta\epsilon\ll1$。离开这个条件，仍套用上式就是把近似用出了适用区间。

大$\Delta'$时改用$\tilde\psi''\sim\tilde\psi/\epsilon^2$，相同两条方程变为
\begin{equation}
 \tilde\phi\sim{K\epsilon\over\hat\gamma}\tilde\psi,\qquad
 \hat\gamma\sim K\epsilon,\qquad
 \hat\gamma\sim{S_a^{-1}\over\epsilon^2}.
\end{equation}
因此
\begin{equation}
 \epsilon\sim S_a^{-1/3}K^{-1/3},\qquad
 \hat\gamma\sim S_a^{-1/3}K^{2/3}.
\end{equation}
这些只是经典电阻内层的量级推导，没有给出匹配积分的数值常数。加入Hall、惯性、FLR或黏性后，内层主阶平衡可能改变，应回到方程重新排序；后面的Hall撕裂和双层匹配专题正是在这个基础上继续。\citep{Baalrud2011,L262}

\subsection{把特征值计算当作推导的延续}
实际求增长率时，先给出平衡$\psi_0(y)$，再选沿片波数$k$，将所有扰动写成$\tilde f(y)e^{ikx+\gamma t}$。导数替换是$\partial_t\to\gamma$、$\partial_x\to ik$，而$\partial_y$仍是待离散的微分算子。边界远离共振面时要求扰动衰减，或采用与具体装置一致的壁面条件。

于是可写成$\mathsf A\vct q=\gamma\mathsf B\vct q$。这里的$\mathsf B$可能包含$1-d_e^2D$及涡量的$D$，不是总能量中的磁场。必须检查边界距离、网格是否分辨内层，以及选到的本征函数是否有撕裂所需的奇偶性。特征值程序返回一个正实部，并不自动证明它是目标撕裂模。
